\scanpage{23}
% L1–3
En fait, le $E$ envisagé a \deleted{nécessairement} une
forme quadratique canonique, car chaque
$B_i\wedge_{\{\pm1\}} k$ a la forme \unsure{standard|standart|bilinéaire} \inserted{de produit scalaire} $(x,y)=xy$.

% L4–5
Je dis qu'on a un isom canonique
\[
\omega_E\simeq\omega_I\wedge\bigwedge_{i\in I}\deleted{\unsure{B_i}} B_i
\]

% L6–7
En effet, par la formule d'associativité on a
\[
\omega_E\simeq\omega_I\wedge\bigwedge_{i\in I}\omega_{E_i}
\]

% L8–9
et d'autre part
\[
\omega_{E_i}\simeq\deleted{\unsure{\omega_{B_i}|\omega_{E_i}}} B_i
\]

% L10–12
% slanted
\margin{faire tomber l'indice $i$}
[ On a en effet, comme $\deleted{\unsure{B}} B_i\subset$ ens des bases de $E_i$,
$B_i\longrightarrow\omega_{E_i}$, et on voit que c'est inj, donc
bijectif \ldots. ]

% L13–14, lower right
\uline{\deleted{Cube combinatoire}} \inserted{épure \deleted{de} $n$-cube}
\uline{orienté}
% Isolated diagonal pen stroke below the text.
\figure{\begin{tikzpicture}\draw (0,0)--(.86,-1.30);\end{tikzpicture}}
% archive sheet number: 22
